\chapter{Tools, Function Calling, and MCP}

The Model Context Protocol (MCP) is a standardised interface for exposing tools, prompts,
and resources to model hosts. It makes the action layer portable: tools can be implemented
once and reused across agents, development environments, desktop applications, and cloud
deployments.

\section{The Function-Calling Loop}

When a model is given access to tools, each reasoning step follows a structured cycle. The
agent framework presents the model with the user request alongside the available tool
schemas. The model determines whether the current state of the task requires a tool
invocation, and if so emits a structured call containing the tool name and a JSON argument
object. The framework validates that call against the schema before passing it to the tool.
The tool runs, and the framework returns its output to the model as an observation. The
model then either synthesises a final answer or initiates another step.

The model does not require knowledge of the tool's implementation. It requires a precise
external contract. This is why schema design, as discussed in the previous chapter, often
has greater influence on agent behaviour than prompt prose: a well-specified schema reduces
ambiguity before the model reasons, whereas a prompt correction must override a confusion
the model has already formed.

\section{MCP Primitives}

The MCP specification defines three principal abstractions.\footnote{MCP specification:
\url{https://modelcontextprotocol.io/specification/2025-06-18}.} Tools are model-invoked
capabilities that interact with external systems, APIs, computations, or stateful actions;
each tool is exposed by a server with a name, metadata, and an input
schema.\footnote{MCP tools specification:
\url{https://modelcontextprotocol.io/specification/2025-06-18/server/tools}.} Resources
are contextual data items exposed by a server and selected by the application or client;
a resource may be a file, a document, a database record, or a repository item. Prompts
are reusable templates and workflow instructions that clients can discover and retrieve
from a server.\footnote{MCP prompts specification:
\url{https://modelcontextprotocol.io/specification/2025-06-18/server/prompts}.}

This distinction maps directly onto the layered architecture of this book. Tools primarily
belong to the action layer. Resources and prompts primarily support the memory and
knowledge layer. All three require security boundaries: because any client that can reach
an MCP server can invoke its tools, the server must enforce authentication and
authorisation independently of the consuming agent.

\section{Local Tools and Remote Tools}

Tools can run inside the agent process or behind a remote endpoint. A local tool is a
plain function co-located with the agent code; it is straightforward to prototype and
debug, but its availability is limited to the agent that defines it and its security
controls are scoped to the process. A remote tool runs behind an API, a serverless
function, an MCP server, or a managed gateway; it introduces network overhead but enables
the tool to be shared across multiple agents, governed through identity and policy
controls, audited centrally, and scaled independently of the agent.

The practical progression is to develop and stabilise a tool locally, establish its
schema and write tests against it, and then promote it to a remote endpoint when it must
be shared or brought under organisational governance.

\section{Validation Before Execution}

Tool call validation occurs in two distinct places, and conflating them is a common
source of production failures. The first is structural validation: the agent framework
checks that the model's emitted arguments conform to the tool's schema before the tool
is invoked. When a tool's input type is declared precisely, using typed function
signatures and constrained field definitions, the framework can enforce this contract
automatically and fail early if the model's output does not conform. The second is
authorisation: an independent component verifies that the caller's identity is permitted
to invoke the requested capability. This separation ensures that schema conformance and
access control are enforced by distinct mechanisms, neither of which depends on the
model's own judgment.

\section{Tool Hallucination}

Tool hallucination occurs when the model invokes the wrong tool, constructs arguments
that the schema does not support, or reasons as though a tool returned information it did
not return. Schema precision is the primary structural mitigation, as established in the
previous chapter. At the observability level, the agent's user interface should surface
which tools were called and with what arguments, making the model's action choices
visible to the user rather than concealing them behind a final answer. At the output
level, the model can be instructed to include a structured attestation field in its
response, for example, a flag indicating that a consequential action such as submitting
a payment was executed, which creates an explicit, checkable link between the model's
stated behaviour and the tool invocations recorded in the trace. These mitigations are
complementary: schema design reduces the rate of hallucination, observability surfaces
instances that occur, and structured attestation makes them detectable in automated
evaluation.

\section{Summary}

MCP standardises the tool interface so capabilities are portable across agents and
deployment environments. Reliable tool use depends on precise schemas, framework-level
validation that is independent of the model, and observability infrastructure that makes
every invocation inspectable. The distinction between local and remote tools is a
governance boundary as much as an architectural one: the point at which a tool moves to
a remote endpoint is the point at which organisational identity, policy, and audit
controls become enforceable.
